\section{Third-Party Shield-Retirement Case Study}
\label{sec:third-party}

The controlled contract in Section~\ref{sec:study} isolates the mechanism.
This section examines the same integrity boundary in an independently
published lifecycle that retires runtime authority.
Carr et al.~\cite{carr_partial_shielding_2023} (AAAI 2023) train a
policy-gradient agent behind a Storm belief-support shield in a partially
observable gridworld.  Their lifecycle later removes the shield either
suddenly or smoothly and compares the released policy with retained-shield
and no-shield controls.  We use the official public implementation and modify
only learner-visible reward records.  The attacker writes
$r'=r+\xi$, with $|\xi|\le\delta$ per record, before the trusted update.
States, observations, shield semantics, dynamics, learner code,
hyperparameters, and evaluation follow the upstream implementation.  The
artifact itemizes every instrumentation and attack patch.  REINFORCE training
runs for 5000 episodes, with retirement at episode 1000.  The TF-Agents
proximal policy optimization (PPO) pipeline retires the shield after $10^5$
environment steps.

\paragraph{Lifecycle fidelity.}
On the obstacle domain ($N{=}6$), the clean condition reproduces the
qualitative ordering of Carr et al.'s Table~1 for REINFORCE: shield retained
0/0 violations during/after training, smooth 326/675, sudden 3277/3784, and
no-shield 3798/3784.  Shield $<$ smooth $\ll$ sudden $\le$ no-shield is the
ordering reported by Carr et al., establishing the lifecycle fidelity of the
instrumented implementation.

\paragraph{Causal isolation at the retirement boundary.}
Under a locked protocol (artifact amendment v1.3), poisoning is confined to
the shielded bootstrap phase.  Reward-record mutation is active only while
the shield enforces the action mask and stops at retirement.  All subsequent
learning is clean.  We evaluate the raw policy at
the switch instant before any unshielded learning (1000 episodes) and again
after the full training budget (5000 episodes).  Each independently trained
clean/poisoned pair is one training-level replicate; configured seeds serve
only as run indices.  Table~\ref{tab:carr-isolation} and
Fig.~\ref{fig:carr-isolation} report the provenance-complete pairs.  A pair is
\emph{positive} when the poisoned at-retirement rate exceeds its clean pair by
at least $0.15$ and the paired McNemar exact test gives $p<0.01$.  All other
outcomes are \emph{non-positive}; \emph{protective} denotes a lower poisoned
rate.

\begin{table}[t]
\centering
\caption{Retirement-boundary causal isolation on the third-party lifecycle
(obstacle, REINFORCE, sudden).  Each row is a provenance-complete paired
training run.  Poisoning stops at retirement, and all later training is
clean.  At-retirement evaluates the raw policy at episode 1000; final
evaluates it at episode 5000.  Rates are violation fractions.}
\label{tab:carr-isolation}
\small
\begin{tabular}{@{}ccccc@{}}
\toprule
 & \multicolumn{2}{c}{at retirement} & \multicolumn{2}{c}{final}\\
\cmidrule(lr){2-3}\cmidrule(lr){4-5}
configured seed & clean & poisoned & clean & poisoned\\
\midrule
1 & 0.216 & 0.777 & 0.236 & 0.235\\
2 & 0.484 & 0.763 & 0.234 & 0.751\\
3$^a$ & 0.241 & 0.242 & 0.747 & 0.228\\
\bottomrule
\end{tabular}
\vspace{1pt}\par\footnotesize $^a$An earlier run-index-3 clean summary
duplicated a separate fix-check record and is excluded.  This row reports a
fresh clean/poisoned pair generated under the locked configuration.  The
archived poisoned realization (0.748 at retirement) lacks a provenance-valid
clean endpoint and is excluded from paired inference.
\end{table}

Two of the three provenance-complete pairs had higher raw-policy violation
rates at retirement under poisoning, by factors of $3.6$ and $1.6$.  The fresh
third pair was unchanged ($0.241\rightarrow0.242$).  Physical violations
remained zero throughout protected execution.  The positive outcomes located
latent policy damage before unshielded learning.  The fresh pair showed that
the fixed attack can fail in a low-disagreement retirement state.  After 4000
episodes of clean post-retirement learning, one positive pair recovered
($0.235$ vs.\ clean $0.236$), one remained worse
($0.751$ vs.\ $0.234$), and the fresh non-positive pair reversed
($0.228$ vs.\ $0.747$).  We denote clean raw-policy/shield disagreement at
retirement by $D_{\mathrm{ret}}$.  In the two positive pairs,
$D_{\mathrm{ret}}$ was $0.025$ and $0.047$ and rose to $0.113$ and $0.109$
under poisoning.  The fresh non-positive pair also had low clean disagreement
($D_{\mathrm{ret}}=0.013$), separating available headroom from attack success.

\begin{figure*}[t]
\centering
\includegraphics[width=\textwidth]{fig6_retirement_isolation_summary.pdf}
\caption{Causal isolation on the third-party lifecycle (REINFORCE, obstacle,
sudden, provenance-complete paired runs).  (a) At-retirement violation fraction of the raw
policy; (b) at-retirement raw-policy/shield disagreement; (c) final
violation fraction after clean post-retirement learning.  Poisoning is
confined to the shielded bootstrap phase.  The run-index-3 bars show the fresh
pair in Table~\ref{tab:carr-isolation}.}
\label{fig:carr-isolation}
\end{figure*}

\paragraph{Learner boundary.}
The same contrast attack ($\delta{=}2$) was evaluated on the third-party PPO
pipeline with action masking and sudden retirement.  It raised post-removal
violations in one of three paired runs, from 1178 to 2439
(McNemar exact $p=2.6\times10^{-149}$).  The first-violation episode changed
from 6 to 0.  The other two runs were unchanged ($p=1.00$ and $0.96$).  In a
separate ten-run shield-on-only isolation
battery, the run indexed by configured seed 101 was positive
($0.219\rightarrow0.754$ at retirement; McNemar exact
$p=3.7\times10^{-113}$), while the other nine entered a high-violation ceiling
regime (clean $0.728$--$0.760$) and were either protective or unchanged.
Figure~\ref{fig:ppo-isolation} reports the full battery.  A ten-run full-phase
PPO battery likewise placed every clean policy near the retirement ceiling
($0.73$--$0.77$) and yielded no positive outcome on the primary retirement
metric.  Its final metric captured a different effect: all ten poisoned
policies remained unsafe ($0.49$--$0.77$), whereas nine clean policies recovered
to approximately $0.22$ under continued learning.

A ten-run shield-on-only battery on off-policy soft actor-critic (SAC)
provided a complementary boundary.  Three runs were positive
($0.207\rightarrow0.561$, $0.218\rightarrow0.513$, and
$0.479\rightarrow0.954$; paired McNemar exact $p<10^{-40}$ for each), three
were protective, and four were unchanged.  All ten clean SAC policies had low
retirement disagreement rather than occupying the ceiling regime.  Four
extension runs generated on a second host were non-positive.  The fixed
attack therefore affected individual PPO and SAC training realizations without
producing reliable success within either family.

\begin{figure*}[t]
\centering
\includegraphics[width=0.85\textwidth]{fig6_ppo_isolation.pdf}
\caption{PPO retirement-boundary isolation on the third-party lifecycle
(shield-on-only contrast poisoning, $\delta{=}2$, paired runs indexed by
run indices 101--110,
sudden retirement).  At-retirement violation fraction of the raw policy for
clean and poisoned runs.  The run indexed by 101 is positive
($0.219\rightarrow0.754$; McNemar exact $p=3.7\times10^{-113}$); the other
nine runs land in the unsafe ceiling regime (clean at-retirement
$0.728$--$0.760$) with no paired headroom, six protective by the paired test and
three unchanged.  The dashed line marks the locked decision boundary
(poisoned at least $0.15$ above clean with paired McNemar $p<0.01$).}
\label{fig:ppo-isolation}
\end{figure*}

\paragraph{Attack budget and state dependence.}
On the full-phase REINFORCE protocol, the effect varied with attack shape,
budget, and training realization.  In the realization indexed by configured
seed~2, $\delta{=}2$ raised post-removal violations from 1127 to 3762
($3.3\times$; McNemar exact $p$ below numerical reporting precision).
Increasing the budget to $\delta{=}10$ produced 3766 violations.  In the
realization indexed by seed~3, $\delta{=}2$ was unchanged (2482 vs.\ 2471;
$p=0.84$), whereas $\delta{=}10$ raised violations to 3700
($p=8.4\times10^{-138}$).  The realization indexed by seed~1 began near the
unsafe ceiling (75.7\%) and showed no added damage at either budget.  Shape
controls were also state-dependent.  The bias variant was protective in the
ceiling realization.  The risk variant was protective in two realizations and
raised violations from 1127 to 2457 in the realization indexed by seed~2.

An independent Linux-host batch with identical dependency versions contained
one low-clean-violation realization and two ceiling realizations.  The former
showed the contrast effect at retirement ($0.232\rightarrow0.501$) and after
removal (1157 $\rightarrow$ 2451; $p=4.6\times10^{-159}$).  The latter two were
unchanged at $\delta{=}2$ and protective at $\delta{=}10$.  The run index of
the low-clean-violation regime differed across hosts.  Repeating the clean
configuration also entered a different trajectory mode ($0.776$ rather than
$0.232$ at retirement).  Dose comparisons are therefore confined to the
locked local batch.  The independent batch shows that attackability follows
the realized retirement state more closely than the nominal seed identifier.

\begin{figure*}[t]
\centering
\includegraphics[width=\textwidth]{fig7_dose_response.pdf}
\caption{Budget- and state-dependent effects on the full-phase REINFORCE
protocol (paired within one locked code version).  (a) Contrast
attack dose response; (b) attack-shape negative controls at $\delta{=}2$.
The run-index-2 realization saturates at $\delta{=}2$; the run-index-3
realization requires $\delta{=}10$; the ceiling realization shows no additional damage at
either budget, and bias/risk controls are protective or strictly weaker.}
\label{fig:dose-response}
\end{figure*}

\paragraph{Escape condition.}
Under the same full-scope poisoning used in the full-phase battery, a retained
shield kept physical violations at zero before and after training in all three
paired runs.  The corresponding sudden-retirement runs exposed final violation
fractions of $0.490$, $0.752$, and $0.744$ after removal.  Resident authority
therefore contained the downstream
consequence of full-scope contamination; removing it exposes the same
consequence as an unsafe released policy.

\paragraph{Interpretation.}
Carr et al.'s workflow supplies the lifecycle component: protected bootstrap
followed by retirement of the shield's action authority.  The controlled
study in Section~\ref{sec:study} supplies the adversarial-provenance
mechanism.  Their composition exposes an integrity boundary at authority
transition.  Carr et al.'s original evaluation concerns shield-removal
behavior under unpoisoned learning; the case study concerns evidence reuse
across the same transition under adversarial update provenance.
Absolute rates differ across instrumented and non-instrumented code versions
because the unmasked-disagreement probe requires an independent belief
tracker; all clean-vs-poisoned comparisons above are paired within a locked
code version and are unaffected by that version change.

Training outcomes also varied across platforms.  In a ten-run REINFORCE
shield-on-only battery on Linux, three low-clean-violation realizations were
positive ($0.240\rightarrow0.533$, $0.234\rightarrow0.744$, and
$0.229\rightarrow0.763$), while seven ceiling realizations showed no added
damage.  By contrast, the fresh pair in Table~\ref{tab:carr-isolation} was
non-positive despite low clean disagreement.  The same configured seed and
code can enter different trajectory modes across hosts or backends.  We
therefore treat each independently trained clean/poisoned pair as a replicate
and use nominal seeds only as run indices.  Episode-level McNemar tests
describe outcomes conditional on a trained pair; the multi-run batteries
provide the training-level replication evidence.

The broader cross-tabulation reinforced this state dependence.  Across 53
locked paired runs spanning both scopes, all three learner families, and both
domains, 12 of 20 runs with clean at-retirement violation below $0.5$ were
positive, whereas none of 33 ceiling runs was positive
(Fig.~\ref{fig:two-by-two}).  The eight low-rate exceptions comprise six SAC
runs and two protective runs in the second environment.  Low clean violation
therefore marks available headroom but does not determine the attack outcome.

The 50 obstacle runs with disagreement instrumentation provided a sharper
retirement-state marker.  Clean $D_{\mathrm{ret}}$ was strongly associated
with clean violation (Spearman $0.976$,
$p=1.3\times10^{-33}$) and negatively associated with attack headroom
(Spearman $-0.547$, $p=3.9\times10^{-5}$).  Every positive outcome had clean
$D_{\mathrm{ret}}\le0.047$, while non-positive outcomes extended to $0.111$.
Under the fixed bands, all 11 positive outcomes fell in the low-disagreement
region ($D_{\mathrm{ret}}<0.06$), and none of 32 high-disagreement runs was
positive.  The area under the receiver operating characteristic curve (AUC)
was $0.953$
(Mann--Whitney $p=5.5\times10^{-6}$).  Family-stratified counts clarify this
statistic.  REINFORCE had seven positive low-band runs and 13 non-positive
high-band runs.  PPO had one positive low-band run and 19 non-positive
high-band runs.  All ten SAC runs were low-band, with three positive and seven
non-positive outcomes.  Thus high disagreement acted as a
strong exclusion condition in this corpus, whereas low disagreement defined
an opportunity region.  Eleven of 18 low-band corpus runs were positive; the
fresh pair, analyzed separately from the 50 disjoint namespaces, gave an
additional low-band non-positive realization.

The protocol chronology bounds the statistical interpretation.  The
$0.06/0.09$ bands were fixed after analysis of 46 existing obstacle rows and
before four SAC extension pairs with run indices 407--410.  All four
extensions were low-band and non-positive.  They strengthened the
non-sufficiency result without prospectively testing high-disagreement
exclusion.  Accordingly, AUC $0.953$ describes discrimination within the
combined locked corpus, not external classifier validation.  Because
$D_{\mathrm{ret}}$ is measured after the shielded bootstrap, it characterizes
the realized retirement state.  An online attack trigger would require
earlier disagreement trajectories that predict the eventual opportunity
region.  The marker is also obstacle-scoped.  Two protective runs in the
second environment had clean disagreement $0.476$ and $0.157$ despite clean
violation below $0.5$.

\begin{figure*}[t]
\centering
\includegraphics[width=0.72\textwidth]{fig9_mechanism_divergence.pdf}
\caption{Clean at-retirement raw-policy/shield disagreement marks an
opportunity region across the obstacle corpus (both scopes, all three learner
families).  Each point is one paired clean/poisoned run; filled markers meet
the positive-outcome rule, and hollow markers do not.
Filled markers all have clean disagreement $\le0.047$; the dashed lines mark
the fixed low/high disagreement bands ($<0.06$ and $\ge0.09$).  No
high-disagreement run is positive, while seven of 18 low-disagreement runs are
non-positive.}
\label{fig:mechanism}
\end{figure*}



\begin{figure*}[t]
\centering
\includegraphics[width=\textwidth]{fig8_2x2_susceptibility.pdf}
\caption{Two-by-two opportunity-region cross-tabulation across all locked
paired runs (both scopes, all three learner families, both domains).  Each
point is one paired clean/poisoned run; filled markers meet the positive rule
(poisoned at-retirement at least $0.15$ above clean at retirement with paired
McNemar exact $p<0.01$), hollow markers do not.  The dashed rule line is
$y=x+0.15$ and the vertical dashed line marks the clean at-retirement $0.5$
boundary: positive outcomes concentrate below $0.5$, and no ceiling run is
positive.}
\label{fig:two-by-two}
\end{figure*}

\paragraph{Environment boundary.}
On a second gridworld of the same third-party suite (surveillance, $N{=}6$,
$R{=}3$), a clean reproduction again satisfied the authors' lifecycle
ordering: shield retained
$0/0$, smooth $0.129$, sudden $0.246$, and no-shield $0.252$ (retained $<$
smooth $\ll$ sudden $\le$ no-shield).  The fixed contrast attack at
$\delta{=}2$ raised post-removal violations in one of three independent paired
runs ($0.000 \rightarrow 0.272$, McNemar exact $p=2.6\times10^{-82}$) and was
protective in the other two ($0.267 \rightarrow 0.000$,
$p=8.4\times10^{-81}$; $0.135 \rightarrow 0.038$,
$p=1.3\times10^{-15}$).  The lifecycle ordering extended to the second
environment, whereas the exploit effect remained environment- and
realization-dependent.
